% Poster: our framework -- Simultaneous Illumination Normalization and
% Sparse Image Compression via Unfolded Recovery
% Dimensions: 48 in (wide) x 36 in (high), landscape.

\documentclass[final]{beamer}

% ====================
% Packages
% ====================

\usepackage[T1]{fontenc}
\usepackage{lmodern}
% 48 in x 36 in  ==  121.92 cm x 91.44 cm  (landscape)
\usepackage[size=custom,width=121.92,height=91.44,scale=1.05]{beamerposter}
\usetheme{thayer}
\usecolortheme{thayer}
\usepackage{graphicx}
\usepackage{booktabs}
\usepackage[numbers]{natbib}
\usepackage{tikz}
\usetikzlibrary{arrows.meta, positioning, calc, shapes.geometric}
\usepackage{pgfplots}
\pgfplotsset{compat=1.14}
\usepackage{anyfontsize}
\usepackage{amsmath}
\usepackage{amssymb}
\usepackage{mathtools}

% Shrink the headline title so it fits on a SINGLE line without
% colliding with the Dartmouth Engineering logo at the top-right
% of the page.
\setbeamerfont{headline title}{size={\fontsize{52}{60}\selectfont},series=\bfseries}
\setbeamerfont{headline author}{size=\large}
\setbeamerfont{headline institute}{size=\small}

% ====================
% Lengths
% ====================

\newlength{\sepwidth}
\newlength{\colwidth}
\setlength{\sepwidth}{0.022\paperwidth}
\setlength{\colwidth}{0.225\paperwidth}

\newcommand{\separatorcolumn}{
  \begin{column}{\sepwidth}
\end{column}}

% ====================
% Title
% ====================

\title{Simultaneous Illumination Normalization and Sparse Image Compression via Unfolded Recovery}

\author{Farhan Sadeek\inst{1}}

\institute[Dartmouth]{\inst{1} Thayer School of Engineering, Dartmouth College, Hanover, NH 03755, USA}

% ====================
% Footer
% ====================

\footercontent{
  ENGS~109 Final Project: Joint Compressed Sensing and Illumination Recovery \hfill
  Dartmouth College, Thayer School of Engineering \hfill
\href{mailto:farhan.sadeek.29@dartmouth.edu}{farhan.sadeek.29@dartmouth.edu}}

% ====================
% Body
% ====================

\begin{document}

\begin{frame}[t]
  \begin{columns}[t]
    \separatorcolumn

    % ============================================================
    % COLUMN 1 -- Motivation, Background, Problem
    % ============================================================
    \begin{column}{\colwidth}

      \begin{block}{Motivation}

        \textbf{Cameras sample then throw away.} A JPEG keeps only
        $5$--$10\%$ of DCT coefficients, and the ISP further reshapes
        the data through demosaicing, tone-mapping, and white-balance.
        Compressed sensing reverses this: $M\!\ll\! N$ random
        measurements suffice if the signal is sparse
        \citep{candes2006stable,donoho2006compressed}.

        \medskip
        \textbf{The sequential pipeline wastes its prior:}
      \begin{enumerate}\setlength{\itemsep}{0.15em}
          \item RAW sensors measure $\boldsymbol{x}=\boldsymbol{s}\odot\boldsymbol{g}$, the
            \emph{product} of sparse reflectance and a smooth
            illumination field.
          \item CS-then-correct enforces $\ell_1$ on $\boldsymbol{x}$, but
            $\boldsymbol{x}$ is \emph{not} DCT-sparse; multiplication by
            $\boldsymbol{g}$ destroys the sparsity the regularizer relies on.
          \item Post-hoc column-mean correction can rescale but
            cannot restore information already shrunken away.
        \end{enumerate}

      \end{block}

      \medskip

      \begin{alertblock}{Our Contributions}

      \begin{enumerate}\setlength{\itemsep}{0.2em}
          \item \textbf{our framework objective.} A single $\ell_1{+}$smooth-prior
            problem in $(\boldsymbol{c},\boldsymbol{g})$ solved by block-coordinate
            alternation: FISTA on $\boldsymbol{c}$, closed-form ridge on $\boldsymbol{g}$.
          \item \textbf{Deep-unfolded inner solver.} A $K{=}10$-layer
            LISTA \citep{gregor2010lista} with Adam, gradient clipping
            and best-checkpoint restore for stability.
          \item \textbf{Scene-dependent benchmark.} Phase transition,
            multi-image rate-distortion, LISTA-vs-classical, and
            per-scene illumination ablation on $5$ natural images.
          \item \textbf{Honest failure-mode analysis.} Joint catastrophically
            fails on near-uniform scenes (moon, $-8$\,dB); we trace
            it to rank-deficiency of $\boldsymbol{B}^{\!\top}\!\boldsymbol{B}$.
        \end{enumerate}

      \end{alertblock}

      \medskip

      \begin{exampleblock}{Big Takeaway}

        \large
        \textbf{Estimating reflectance and illumination
          \emph{jointly} beats the conventional CS-then-correct
        pipeline by up to $+4.7$\,dB on textured scenes ---
        \emph{when} the scene has enough column-wise energy to
        identify the gain.}

        \medskip\normalsize
      \begin{itemize}\setlength{\itemsep}{0.15em}
          \item \textbf{Textured wins:} $+1.1$ to $+4.7$\,dB (cameraman, astronaut).
          \item \textbf{LISTA speedup:} $K{=}10$ matches FISTA at $22$ iters.
          \item \textbf{Failure mode:} $-8$\,dB on near-uniform (moon)
            -- needs an identifiability gate.
        \end{itemize}

      \end{exampleblock}

      \medskip

      \begin{block}{Background: Compressed Sensing}

        With $\boldsymbol{x}{=}\boldsymbol{\Psi}\boldsymbol{c}$ and $\|\boldsymbol{c}\|_0\!\leq\!S$, the
        forward model is
        $\boldsymbol{y}{=}\boldsymbol{A}\boldsymbol{x}{+}\boldsymbol{w}$, $\boldsymbol{A}{\in}\mathbb R^{M{\times}N}$.
        If $\boldsymbol{A}$ satisfies the RIP with
        $\delta_{2S}\!<\!\sqrt 2{-}1$, the Lasso relaxation
        \[
          \boldsymbol{c}_\lambda
          =\arg\min_{\boldsymbol{c}}\;\tfrac12\|\boldsymbol{A}\boldsymbol{\Psi}\boldsymbol{c}-\boldsymbol{y}\|_2^2
          +\lambda\|\boldsymbol{c}\|_1
        \]
        recovers $\boldsymbol{c}$ to within $O(\epsilon)$
        \citep{candes2006stable}. We use the 2D DCT as $\boldsymbol{\Psi}$, a
        strong natural-image sparsifying basis.

      \end{block}

      \medskip

      \begin{block}{The Sequential-Pipeline Problem}

        A conventional pipeline recovers $\boldsymbol{x}$ first, then divides
        out an illumination estimate. Two errors compound:
      \begin{itemize}\setlength{\itemsep}{0.12em}
          \item \textbf{Sparsity broken at the prior.} Pointwise
            multiplication by $\boldsymbol{g}$ kills DCT-domain sparsity, so
            $\ell_1$ over-shrinks high-illumination coefficients.
          \item \textbf{Correction is shape-rigid.} A column-mean
            divisor copies the gradient but cannot restore the
            information the over-shrinkage has already destroyed.
        \end{itemize}
        our framework estimates $\boldsymbol{c}$ and $\boldsymbol{g}$ together so the
        $\ell_1$ prior sees the true sparse signal $\boldsymbol{s}$.

      \end{block}

    \end{column}

    \separatorcolumn

    % ============================================================
    % COLUMN 2 -- Notation + Method 1 + Method 2
    % ============================================================
    \begin{column}{\colwidth}

      \begin{block}{Notation}

        \begin{tabular}{@{}l@{\hspace{0.9em}}l@{}}
          $\boldsymbol{y}\!\in\!\mathbb R^M$  & CS measurements ($M{\ll}N$) \\
          $\boldsymbol{A}\!\in\!\mathbb R^{M{\times}N}$ & Gaussian sensing matrix \\
          $\boldsymbol{\Psi}$ & 2D DCT sparsifying basis \\
          $\boldsymbol{c}$ & sparse DCT coefficients \\
          $\boldsymbol{s}{=}\boldsymbol{\Psi}\boldsymbol{c}$ & sparse reflectance image \\
          $\boldsymbol{g}\!\in\!\mathbb R^{\sqrt N}$ & per-column illumination gain \\
          $\boldsymbol{x}{=}\boldsymbol{s}\!\odot\!\boldsymbol{g}_\textrm{pix}$ & RAW signal \\
        \end{tabular}

      \end{block}

      \medskip

      \begin{block}{Method 1: our framework Objective}

        \textbf{Replace} the two-stage pipeline with a single objective:
        \[
          \boxed{\;
            \min_{\boldsymbol{c},\boldsymbol{g}}\;
            \tfrac12\big\|\boldsymbol{A}\,\mathrm{diag}(\boldsymbol{g}_\textrm{pix})
            \boldsymbol{\Psi}\boldsymbol{c}-\boldsymbol{y}\big\|_2^2
            +\lambda_c\|\boldsymbol{c}\|_1
            +\lambda_g\|\nabla\boldsymbol{g}\|_2^2
          \;}
        \]
        Jointly non-convex but \textbf{block-convex} in each variable.

        \medskip
        \textbf{$\boldsymbol{c}$-step.} Fix $\boldsymbol{g}$; effective sensing matrix
        $\boldsymbol{A}_{\!\textrm{eff}}{=}\boldsymbol{A}\,\mathrm{diag}(\boldsymbol{g}_\textrm{pix})\boldsymbol{\Psi}$,
        update is a standard Lasso solved by FISTA.

        \medskip
        \textbf{$\boldsymbol{g}$-step.} Fix $\boldsymbol{c}$; with $\hat{\boldsymbol{s}}{=}\boldsymbol{\Psi}\boldsymbol{c}$
        the model is linear in $\boldsymbol{g}$,
        $\boldsymbol{y}{=}\boldsymbol{B}(\hat{\boldsymbol{s}})\boldsymbol{g}{+}\boldsymbol{w}$, and the
        update is closed-form ridge:
        \[
          \boldsymbol{g}^\star
          =\bigl(\boldsymbol{B}^{\!\top}\!\boldsymbol{B}+\lambda_g \boldsymbol{L}\bigr)^{-1}\boldsymbol{B}^{\!\top}\boldsymbol{y},
        \]
        with $\boldsymbol{L}{=}\nabla^{\!\top}\!\nabla$ the 1D discrete
        Laplacian. We clamp $\boldsymbol{g}{\geq}\epsilon_g$ to avoid
        divide-by-zero.

        \medskip
        \textbf{Gain ambiguity.} The model has scalar invariance
        $(\boldsymbol{s},\boldsymbol{g})\!\to\!(a\boldsymbol{s},\boldsymbol{g}/a)$; we resolve at
        evaluation by the optimal positive gain
        $a^\star{=}\langle\hat{\boldsymbol{s}},\boldsymbol{s}_\textrm{true}\rangle/\|\hat{\boldsymbol{s}}\|^2$,
        applied identically to baselines.

      \end{block}

      \medskip

      \begin{block}{Method 2: Deep-Unfolded LISTA Inner Solver}

        Replace the FISTA $\boldsymbol{c}$-step with $K$ learned layers
        \citep{gregor2010lista}:
        \[
          \boldsymbol{c}^{(k+1)}
          =\mathcal S_{\theta_k}\!\bigl(
          \boldsymbol{W}_t\boldsymbol{c}^{(k)}+\boldsymbol{W}_e\boldsymbol{y}\bigr),
          \quad k=0,\dots,K{-}1.
        \]
        Initialise
        $\boldsymbol{W}_e{=}\tfrac1L(\boldsymbol{A}\boldsymbol{\Psi})^{\!\top}$,
        $\boldsymbol{W}_t{=}\boldsymbol{I}{-}\tfrac1L(\boldsymbol{A}\boldsymbol{\Psi})^{\!\top}\boldsymbol{A}\boldsymbol{\Psi}$,
        $\theta_k{=}\lambda/L$ to reproduce ISTA at step zero, then
        train end-to-end on $(\boldsymbol{y},\boldsymbol{c})$ pairs.

        \medskip
        \begin{figure}
          \centering
          \begin{tikzpicture}[
              scale=1.0,
              every node/.style={font=\footnotesize},
              lstep/.style={rounded corners=3pt, draw=black!70, very thick,
                fill=blue!10, minimum width=2.0cm, minimum height=0.9cm,
              font=\small\bfseries, align=center},
              thr/.style={rounded corners=3pt, draw=orange!75!black, very thick,
                fill=orange!20, minimum width=1.5cm, minimum height=0.9cm,
              font=\small\bfseries, align=center},
              flowarr/.style={-{Stealth[length=2.5mm]}, thick, draw=black!65}
            ]
            \node[lstep] (l1) at (0,0)  {LISTA\\layer 1};
            \node[thr]   (s1) at (2.2,0) {$\mathcal S_{\theta_1}$};
            \node[lstep] (l2) at (4.6,0) {LISTA\\layer 2};
            \node[thr]   (s2) at (6.8,0) {$\mathcal S_{\theta_2}$};
            \node          at (8.6,0)   {$\cdots$};
            \node[lstep] (lk) at (10.6,0){LISTA\\layer $K$};
            \node[thr]   (sk) at (12.8,0){$\mathcal S_{\theta_K}$};

            \node[font=\small] (y) at (-2.0,0) {$\boldsymbol{y}$};
            \node[font=\small] (c) at (14.8,0) {$\hat{\boldsymbol{c}}$};

            \draw[flowarr] (y)--(l1);
            \draw[flowarr] (l1)--(s1);
            \draw[flowarr] (s1)--(l2);
            \draw[flowarr] (l2)--(s2);
            \draw[flowarr] (s2)--(8.0,0);
            \draw[flowarr] (9.2,0)--(lk);
            \draw[flowarr] (lk)--(sk);
            \draw[flowarr] (sk)--(c);

            \node[font=\scriptsize, text=black!65, anchor=north]
              at (6.4,-0.9) {tied $\boldsymbol{W}_e,\boldsymbol{W}_t$; per-layer learnable $\theta_k$};
          \end{tikzpicture}
          \caption{$K$-layer LISTA with tied weights and per-layer
            thresholds. Each $\theta_k$ adapts to that layer's effective
            noise; the recurrence reproduces ISTA at initialisation.}
        \end{figure}

        \heading{Stability tricks}

      \begin{itemize}\setlength{\itemsep}{0.12em}
          \item Naive Adam@$5{\times}10^{-3}$: clean descent to
            $0.12$ NMSE, then blow-up to $\sim\!10^3$.
          \item Diagnosis: $\rho(\boldsymbol{W}_t){>}1$ along some
            eigendirection causes the recurrence to diverge.
          \item Fix: (i) lr$\to\!10^{-3}$, (ii) clip-norm $1.0$,
            (iii) restore best-val checkpoint.
        \end{itemize}

      \end{block}

    \end{column}

    \separatorcolumn

    % ============================================================
    % COLUMN 3 -- Pipeline + Algorithm + Phase Transition
    % ============================================================
    \begin{column}{\colwidth}

      \begin{block}{Encoder $\to$ Decoder Pipeline}

        \begin{figure}
          \centering
          \begin{tikzpicture}[
              xscale=1.18,
              every node/.style={font=\footnotesize},
              raw/.style={rounded corners=2pt, draw=black!60, fill=blue!10,
                draw=blue!55!black,
                minimum width=2.8cm, minimum height=0.8cm,
              font=\footnotesize\bfseries, align=center},
              meas/.style={rounded corners=2pt, draw=green!55!black, fill=green!10,
                minimum width=2.8cm, minimum height=0.8cm,
              font=\footnotesize\bfseries, align=center},
              joint/.style={rounded corners=4pt, draw=black!80, very thick,
                fill=yellow!12, minimum width=2.8cm, minimum height=1.4cm,
              font=\footnotesize\bfseries, align=center},
              dec/.style={rounded corners=4pt, draw=red!70!black, very thick,
                fill=red!8, minimum width=2.8cm, minimum height=0.8cm,
              font=\footnotesize\bfseries, align=center},
              flowarr/.style={-{Stealth[length=2.5mm]}, thick, draw=black!65}
            ]
            \node[raw]   (s)   at (0, 1.2) {reflectance $\boldsymbol{s}$};
            \node[raw]   (g)   at (0,-1.2) {illumination $\boldsymbol{g}$};
            \node[meas]  (y)   at (3.4, 0.0) {$\boldsymbol{y}{=}\boldsymbol{A}\boldsymbol{x}{+}\boldsymbol{w}$};
            \node[joint] (alt) at (7.4, 0.0) {block-coord\\$\boldsymbol{c}\!\leftrightarrow\!\boldsymbol{g}$};
            \node[raw]   (sh)  at (11.0, 1.2) {$\hat{\boldsymbol{s}}{=}\boldsymbol{\Psi}\hat{\boldsymbol{c}}$};
            \node[raw]   (gh)  at (11.0,-1.2) {$\hat{\boldsymbol{g}}$};
            \node[dec]   (out) at (14.6, 0.0) {gain-match\\+ PSNR};

            \draw[flowarr] (s.east) -- (y.west);
            \draw[flowarr] (g.east) -- (y.west);
            \draw[flowarr] (y) -- (alt);
            \draw[flowarr] (alt.east) -- (sh.west);
            \draw[flowarr] (alt.east) -- (gh.west);
            \draw[flowarr] (sh) -- (out);
            \draw[flowarr] (gh) -- (out);

            \node[font=\scriptsize, anchor=north, text=black!65] at (1.7,-2.1)
            {$\boldsymbol{x}{=}\boldsymbol{s}\!\odot\!\boldsymbol{g}$};
            \node[font=\scriptsize, anchor=north, text=black!65] at (5.4,-2.1)
            {Gaussian $\boldsymbol{A}$, AWGN};
            \node[font=\scriptsize, anchor=north, text=black!65] at (9.2,-2.1)
            {LISTA / FISTA + ridge};
            \node[font=\scriptsize, anchor=north, text=black!65] at (12.8,-2.1)
            {$a^\star\!=\!\langle\hat{\boldsymbol{s}},\boldsymbol{s}\rangle/\|\hat{\boldsymbol{s}}\|^2$};
          \end{tikzpicture}
          \caption{End-to-end our framework. The RAW signal
            $\boldsymbol{x}{=}\boldsymbol{s}\!\odot\!\boldsymbol{g}$ is measured via Gaussian
            $\boldsymbol{A}$; block-coordinate alternation recovers
            $(\boldsymbol{c},\boldsymbol{g})$ simultaneously; the scalar gain
            ambiguity is resolved at evaluation.}
        \end{figure}

      \end{block}

      \medskip

      \begin{block}{Block-Coordinate Recipe}

      \begin{enumerate}\setlength{\itemsep}{0.18em}
          \item Initialise $\boldsymbol{g}^{(0)}{=}\boldsymbol{1}$,
            $\boldsymbol{c}^{(0)}{=}\boldsymbol{0}$.
          \item \textbf{$\boldsymbol{c}$-update.} Hold $\boldsymbol{g}^{(t)}$; run
            $T_c$ FISTA (or $K$ LISTA) steps on
            $\boldsymbol{A}_{\!\textrm{eff}}(\boldsymbol{g}^{(t)})$.
          \item \textbf{$\boldsymbol{g}$-update.} Hold $\boldsymbol{c}^{(t{+}1)}$; solve
            $(\boldsymbol{B}^{\!\top}\!\boldsymbol{B}{+}\lambda_g\boldsymbol{L})\boldsymbol{g}{=}\boldsymbol{B}^{\!\top}\!\boldsymbol{y}$
            in closed form; clamp to $[\epsilon_g,\infty)$.
          \item Repeat $T_\textrm{outer}{=}8$ outer iterations.
          \item Resolve scalar gain ambiguity at evaluation.
        \end{enumerate}

      \end{block}

      \medskip

      \begin{exampleblock}{Phase Transition (Sec.~3.2)}

        Empirical Donoho--Tanner boundary on synthetic Gaussian CS,
        $N{=}200$, $20$ trials/cell, NMSE${<}10^{-3}$.

        \medskip
        \centering\small
        \renewcommand{\arraystretch}{1.15}
        \begin{tabular}{@{}l r r r r@{}}
          \toprule
          \textbf{$\delta=M/N$} & $0.30$ & $0.50$ & $0.70$ & $0.90$ \\
          \midrule
          OMC $\rho^\star$   & $\mathbf{0.35}$ & $\mathbf{0.40}$ & $0.40$         & $0.50$ \\
          FISTA $\rho^\star$ & $0.20$ & $0.30$ & $\mathbf{0.50}$ & $\mathbf{0.65}$ \\
          \bottomrule
        \end{tabular}

        \medskip
      \begin{itemize}\setlength{\itemsep}{0.12em}
          \item \textbf{OMP wins at low $\delta$:} exact LS-on-support
            beats $\ell_1$ when measurements are scarce.
          \item \textbf{FISTA wins at high $\delta$:} $\ell_1$ keeps
            expanding the feasible region; OMP plateaus on dense
            supports.
          \item \textbf{Crossover} near $\delta{\approx}0.65$.
        \end{itemize}

      \end{exampleblock}

      \medskip

      \begin{block}{Rate-Distortion (Sec.~3.3)}

        $5$-image natural-scene set, SNR${=}30$\,dB, $16{\times}16$
        blocks. Mean PSNR$\pm$std across scenes.

        \medskip
        \centering\small
        \renewcommand{\arraystretch}{1.15}
        \begin{tabular}{@{}l r r r r@{}}
          \toprule
          $\delta$ & $0.20$ & $0.40$ & $0.50$ & $0.70$ \\
          \midrule
          OMP   & $17.4{\pm}4.3$ & $20.4{\pm}3.7$ & $21.7{\pm}3.2$ & $24.9{\pm}3.0$ \\
          FISTA & $\mathbf{19.4{\pm}4.5}$ & $\mathbf{23.6{\pm}4.1}$ & $\mathbf{25.1{\pm}3.9}$ & $\mathbf{27.6{\pm}3.5}$ \\
          ADMM  & $17.6{\pm}4.8$ & $23.5{\pm}4.4$ & $25.0{\pm}4.0$ & $27.5{\pm}3.7$ \\
          \bottomrule
        \end{tabular}

        \medskip
        \raggedright
        \textbf{Verdict.} FISTA and ADMM agree within $0.5$\,dB above
        $\delta{=}0.2$ (same Lasso objective); both beat OMP by
        $2$--$3$\,dB.  Wide std reflects scene difficulty: moon is
        easy ($>30$\,dB), coins is hard ($\sim20$\,dB).

      \end{block}

    \end{column}

    \separatorcolumn

    % ============================================================
    % COLUMN 4 -- LISTA + Joint vs Sequential + Limitations + Refs
    % ============================================================
    \begin{column}{\colwidth}

      \begin{block}{LISTA vs.\ Classical at Matched Compute}

      \begin{itemize}\setlength{\itemsep}{0.12em}
          \item \textbf{Setup.} $K{=}10$ tied-weight LISTA, $5000$
            synthetic $(\boldsymbol{y},\boldsymbol{c})$ pairs, $N{=}200, M{=}80, S{=}10$,
            SNR${=}30$\,dB.
        \end{itemize}

        \medskip
        \centering\small
        \renewcommand{\arraystretch}{1.18}
        \begin{tabular}{@{}l c c@{}}
          \toprule
          \textbf{Solver} & \textbf{Iters/Layers} & \textbf{Val NMSE} \\
          \midrule
          ISTA   & $10$  & $0.467$ \\
          FISTA  & $10$  & $0.337$ \\
          \textbf{LISTA}  & $10$  & $\mathbf{0.050}$ \\
          \midrule
          FISTA (matched) & $22$  & $\approx 0.05$ \\
          FISTA (conv.)   & $500$ & $0.005$ \\
          \bottomrule
        \end{tabular}

        \medskip
        \raggedright
        \textbf{$2.2\times$ inference speedup} at matched recovery;
        LISTA still trails the converged $\ell_1$ minimum by $10\times$,
        but amortises that gap into a $10\times$ smaller cost at
        deployment.

      \end{block}

      \medskip

      \begin{exampleblock}{Joint vs.\ Sequential: Per-Scene PSNR Gain (dB)}

        $8{\times}$ horizontal illumination gradient, SNR${=}25$\,dB,
        $16{\times}16$ center patch per scene.
        $\Delta{=}\textrm{PSNR}_\textrm{joint}{-}\textrm{PSNR}_\textrm{seq}$.

        \medskip
        \begin{figure}
          \centering
          \begin{tikzpicture}
            \begin{axis}[
                width=0.95\textwidth,
                height=0.45\textwidth,
                ybar,
                bar width=4.5pt,
                enlarge x limits=0.10,
                ymin=-11, ymax=6,
                ytick={-10,-8,-6,-4,-2,0,2,4},
                ylabel={$\Delta$ PSNR (dB)},
                xlabel={Measurement rate $\delta$},
                symbolic x coords={0.20,0.30,0.40,0.50,0.60},
                xtick=data,
                legend style={at={(0.02,0.02)}, anchor=south west,
                  legend columns=3, font=\scriptsize,
                  draw=black!40, fill=white, fill opacity=0.9,
                text opacity=1, rounded corners=2pt},
                legend cell align=left,
                ymajorgrids=true,
                grid style={dashed,gray!30},
                tick label style={font=\normalsize},
                label style={font=\normalsize},
                axis line style={black!60},
              ]
              \addplot[fill=blue!55,draw=blue!70!black] coordinates {
                (0.20,2.87) (0.30,4.65) (0.40,1.13) (0.50,1.08) (0.60,3.07)
              };
              \addplot[fill=orange!60,draw=orange!70!black] coordinates {
                (0.20,0.54) (0.30,0.18) (0.40,1.84) (0.50,2.31) (0.60,2.21)
              };
              \addplot[fill=green!50,draw=green!55!black] coordinates {
                (0.20,-1.02) (0.30,-1.73) (0.40,0.39) (0.50,0.74) (0.60,0.32)
              };
              \addplot[fill=purple!40,draw=purple!60!black] coordinates {
                (0.20,-0.79) (0.30,-0.48) (0.40,-0.23) (0.50,0.74) (0.60,0.43)
              };
              \addplot[fill=red!55,draw=red!70!black] coordinates {
                (0.20,-10.03) (0.30,-9.51) (0.40,-7.76) (0.50,-9.05) (0.60,-8.37)
              };
              \draw[black!50, dashed, very thick]
                ({rel axis cs:0,0} -| {axis cs:0.20,0}) --
                ({rel axis cs:1,0} -| {axis cs:0.60,0});
              \legend{cameraman,astronaut,coins,page,moon}
            \end{axis}
          \end{tikzpicture}
          \caption{Per-scene PSNR gain of our framework over the sequential
            baseline. Textured scenes (cameraman, astronaut) gain
            $+1.1$ to $+4.7$\,dB; coins/page are roughly a wash;
            \textbf{moon collapses by $\sim 8$\,dB} -- a real,
            mechanism-explainable failure mode.}
        \end{figure}

      \end{exampleblock}

      \medskip

      \begin{block}{Why Moon Fails (and Why That's Informative)}

        Moon's centre patch is mostly black sky with a thin sliver of
        bright surface. After $8{\times}$ gradient multiplication, most
        columns carry near-zero energy, so
        $\boldsymbol{B}^{\!\top}\!\boldsymbol{B}$ is severely rank-deficient and
        per-column gains are non-identifiable. The alternation then
        drives $\boldsymbol{g}$ to a high-variance solution that fits noise.
        The sequential pipeline avoids this by \emph{never} estimating
        $\boldsymbol{g}$ from data -- its column-mean divisor is essentially
        a copy of the gradient. \textbf{When the scene cannot identify
        $\boldsymbol{g}$, refusing to estimate it is a virtue.}

        \medskip
        \textbf{Deployment rule:} gate the joint solution on
        $\kappa(\boldsymbol{B}^{\!\top}\!\boldsymbol{B}){<}\tau$; fall back to
        sequential otherwise. Not reported in prior CS$+$ISP literature.

      \end{block}

      \medskip

      \begin{block}{Limitations \& Future Work}

      \begin{itemize}\setlength{\itemsep}{0.12em}
          \item \textbf{Scene-dependent failure} on near-uniform images
            (moon $-8$\,dB) needs an identifiability gate.
          \item \textbf{1D illumination model} suits horizontal
            gradients; 2D vignettes need a richer $\boldsymbol{g}$ prior.
          \item \textbf{Synthetic-trained LISTA} would need
            distribution-wide training for noise generalisation.
          \item \textbf{Simulated measurements only}; physical
            coded-aperture validation remains future work.
          \item \textbf{Outer-loop unfolding} (not just $\boldsymbol{c}$-step)
            would yield spatially-adaptive thresholds.
        \end{itemize}

      \end{block}

      \medskip

      \begin{block}{References}

        \nocite{*}
        \renewcommand{\bibsep}{1pt}
        {\fontsize{14pt}{17pt}\selectfont
        \bibliographystyle{plainnat}\bibliography{poster}}

      \end{block}

    \end{column}

    \separatorcolumn
  \end{columns}
\end{frame}

\end{document}
